\centerline{\bf \Huge Полюсы и поляры 2}


\problem{1}
\textbf{Теорема Брокара.}
Четырёхугольник $A B C D$ вписан в окружность с центром $O$. Пусть $P=A B \cap C D$, $Q=A D \cap B C, R=A C \cap B D$.

(a) Докажите, что поляры точек $P, Q$ и $R$ есть прямые $Q R, P R$ и $P Q$ соответственно.

(b) Докажите, что $O$ является ортоцентром треугольника $P Q R$.

\problem{2}
Пусть в четырёхугольник $A B C D$ вписана окружность $\omega$, которая касается сторон $A B, B C, C D$ и $D A$ в точках $K, L, M$ и $N$ соответственно. Прямые $K L$ и $M N$ пересекаются в точке $S$, а прямые $L M$ и $N K$ - в точке $T$. Прямые $A B$ и $C D$ пересекаются в точке $P, B C$ и $A D$ - в точке $Q$. Докажите, что тогда

(a) $P, Q, S$ и $T$ лежат на одной прямой;

(b) прямые $A C, B D, K M$ и $L N$ пересекаются в одной точке.

\problem{3}
Прямая $\ell$ касается вписанной окружности ($I$ - её центр) треугольника $A B C$. $A_0$ лежит на прямой $\ell$, так что $\angle A I A_0=90^{\circ}$. Аналогично, определим точки $B_0, C_0$. Докажите, что прямые $A A_0, B B_0, C C_0$ пересекаются в одной точке.

\problem{4}
В треугольнике $A B C$ вписанная окружность $\omega$ касается сторон $A B$ и $A C$ в точках $C_1$ и $B_1$. Оказалось, что прямые $B C, B_1 C_1$ и касательная к $A B C$ в точке $A$ пересекаются в одной точке $D$. Докажите, что $D$ лежит на прямой, соединяющей центры $\omega$ и $(A B C)$ .

Отношение «полюс-поляра» задаёт биекцию между точками и прямыми на проективной плоскости. Замена точек на соответствующие им прямые и наоборот называется полярным преобразованием. При полярном преобразовании одни утверждения переходят в другие. Классическим примером является теоремы Паскаля и Брианшона.

\problem{5}
Осознайте, что теоремы
\href{https://ru.wikipedia.org/wiki/%D0%A2%D0%B5%D0%BE%D1%80%D0%B5%D0%BC%D0%B0_%D0%9F%D0%B0%D1%81%D0%BA%D0%B0%D0%BB%D1%8F}{Паскаля}
и
\href{https://ru.wikipedia.org/wiki/%D0%A2%D0%B5%D0%BE%D1%80%D0%B5%D0%BC%D0%B0_%D0%91%D1%80%D0%B8%D0%B0%D0%BD%D1%88%D0%BE%D0%BD%D0%B0}{Брианшона}
переходят друг в друга при полярном преобразовании.


\problem{6} 
(а)
Докажите, что полярное преобразование сохраняет углы между прямыми. Более конкретно, если $a$ и $b$ --- поляры точек $A$ и $B$, то $\angle(O A, O B)= \angle(a, b)$, где $O$ центр окружности. Это позволяет следить за углами при полярном преобразовании.

\noindent
(б) 
Докажите, что полярное преобразование сохраняет двойные отношения.